\documentclass[12pt]{extarticle}
\def\activityid{F03-M-v2}\usepackage[letterpaper,margin=.65in,top=.60in,bottom=.65in]{geometry}
\usepackage[T1]{fontenc}\usepackage[default]{sourcesanspro}
\usepackage{microtype,tikz,array,tabularx,fancyhdr,amsmath,amssymb}
\usepackage[hidelinks]{hyperref}\usetikzlibrary{calc,arrows.meta}
\setlength{\parindent}{0pt}\setlength{\parskip}{7pt}\setlength{\headheight}{0pt}\setlength{\footskip}{26pt}
\setlength{\emergencystretch}{2em}\pagestyle{fancy}\fancyhf{}\renewcommand{\headrulewidth}{0pt}\renewcommand{\footrulewidth}{.4pt}
\fancyfoot[L]{\scriptsize Bellingham Math Circle / Week 3 / \activityid}\fancyfoot[R]{\scriptsize \thepage}
\definecolor{ink}{gray}{.12}\definecolor{soft}{gray}{.95}\color{ink}
\newcommand{\PageTitle}[2]{{\fontsize{10}{12}\selectfont\bfseries WEEK 3\quad CODE MACHINES\hfill #1\par}\vspace{3pt}{\fontsize{26}{29}\selectfont\bfseries #2\par}\vspace{2pt}}
\newcommand{\NameLine}{{\small Name: \rule{2.65in}{.4pt}\hfill Date: \rule{1.35in}{.4pt}\par}\vspace{4pt}}
\newcommand{\Task}[2]{\par\vspace{3pt}{\large\bfseries #1}\par #2\par}
\newcommand{\Subhead}[1]{\par\vspace{3pt}{\large\bfseries #1}\par}
\newcommand{\RuleBox}[1]{\par\begingroup\setlength{\fboxsep}{8pt}\colorbox{soft}{\parbox{\dimexpr\linewidth-16pt\relax}{#1}}\endgroup\par}
\newcommand{\WriteBox}[2]{\par\textbf{#1}\par\vspace{2pt}\begin{tikzpicture}\draw[black!40,line width=.5pt] (0,0) rectangle (\dimexpr\linewidth-.5pt\relax,#2);\end{tikzpicture}\par}
\newcommand{\StudentFooter}{\vfill{\small Try it with objects. Explain by pointing, talking, or drawing.}\par}
\newcommand{\LampRing}[3]{\begin{tikzpicture}[x=1in,y=1in]
\foreach \i in {1,...,#1}{\coordinate (v\i) at ({90-360*(\i-1)/#1}:#2);}
\foreach \i in {1,...,#1}{\pgfmathtruncatemacro{\j}{mod(\i,#1)+1}\draw[thick] (v\i)--node[midway,fill=white,inner sep=2pt,font=\small]{\i\j}(v\j);}
\foreach \i in {1,...,#1}{\node[circle,draw,fill=white,minimum size=.52in,inner sep=0pt,font=\bfseries] at(v\i){\i};}
\foreach \i in {#3}{\node[circle,draw,fill=ink,text=white,minimum size=.52in,inner sep=0pt,font=\bfseries] at(v\i){\i};}
\end{tikzpicture}}
\newcommand{\Lights}[1]{\begin{tikzpicture}[x=.55in,y=.55in,baseline=-.10in]\foreach \v [count=\i] in {#1}{\ifnum\v=1\fill (\i,0) circle(.16in);\else\draw (\i,0) circle(.16in);\fi}\end{tikzpicture}}
\newcommand{\ClockRing}[2]{\begin{tikzpicture}[x=1in,y=1in]\draw[black!30] (0,0) circle(#2);\pgfmathtruncatemacro{\n}{#1-1}\foreach\i in {0,...,\n}{\node[circle,draw,fill=white,minimum size=.35in,inner sep=0pt] at({90-360*\i/#1}:#2){\i};}\draw[-{Stealth},thick] (-.35,.15) arc(155:25:.38);\node[font=\small] at(0,-.18){clockwise};\end{tikzpicture}}
\newcommand{\Key}[2]{\begin{center}\renewcommand{\arraystretch}{1.55}\begin{tabular}{c|*{#1}{c}}in&#2\end{tabular}\end{center}}

\begin{document}

\PageTitle{GRADES 2--3 / 1}{A code you can undo}\NameLine
\RuleBox{Use cards A, B, C, D. A key uses every letter exactly once in its output row. One turn replaces \textbf{every} letter at once. Keep spaces and letter positions fixed.}
\begin{center}\renewcommand{\arraystretch}{1.6}\begin{tabular}{c|cccc}Input&A&B&C&D\\\hline Key P&B&C&A&D\\Key Q&B&A&D&C\end{tabular}\end{center}
\Task{1. What stays the same?}{Encode ABBA with P. Could any valid key turn ABBA into ABCD? Explain. Then decode P's output CAAC.}
\WriteBox{My messages and a reason:}{.9in}
\Task{2. Keep coding the code.}{Start with ABCD. Apply P again and again until \textbf{all four letters} are back in their starting places. Do the same with Q. Count the first positive return.}
\begin{center}\renewcommand{\arraystretch}{1.7}\begin{tabular}{c|p{2.4in}|p{2.4in}}Turn&Key P&Key Q\\\hline 0&ABCD&ABCD\\1&&\\2&&\\3&&\\4&&\end{tabular}\end{center}
\textbf{Look closer:} Follow just A through P. What happens to D? Is one letter returning enough to say the \emph{whole machine} has reset?
\StudentFooter
\newpage
\PageTitle{GRADES 2--3 / 2}{Build a slow machine}
\Task{3. Make a key that takes four turns.}{Write each of A, B, C, D exactly once under the line. Test it with the whole message ABCD. Draw arrows to show the route of each letter.}
\begin{center}\renewcommand{\arraystretch}{1.8}\begin{tabular}{c|p{.6in}p{.6in}p{.6in}p{.6in}}Input&A&B&C&D\\\hline My key&&&&\end{tabular}\end{center}
\WriteBox{Arrow loops and my test:}{1.1in}
\Task{4. Can four letters take six turns?}{Try it. Then sort your keys by their loops. A letter that stays put is a loop of size 1. What are all the ways to split four letters into loops?}
\begin{center}\renewcommand{\arraystretch}{1.65}\begin{tabular}{p{2.4in}|p{2.4in}}Loop sizes&First turn when all are home\\\hline&\vspace{0pt}\\&\vspace{0pt}\\&\vspace{0pt}\\&\vspace{0pt}\\&\vspace{0pt}\\\end{tabular}\end{center}
\textbf{Go further:} Add a fifth letter. Can you build a machine that really does take six turns? Try two separate loops.
\StudentFooter

\end{document}
